\section{Conclusion and Future Work}
\label{sec:conclusion}

We have presented $\protocol$, a verifiable computation protocol for matrix
multiplication that achieves a SNARK circuit size of $\mathcal{O}(tk)$, where
$t$ is a small constant determined by the security parameter and $k$ is the
matrix dimension.
The protocol combines Freivalds' randomized reduction with proximity testing
over linear error-correcting codes, cleanly decoupling the expensive matrix
arithmetic (performed outside the SNARK by the prover) from the lightweight
algebraic checks (arithmetized inside the SNARK).
We proved that $\protocol$ achieves perfect completeness and computational
soundness with error
$\epsilon_{\mathsf{bind}} + (k-1)/|\mathbb{F}| + 4(1-\delta)^t$.

Our experimental evaluation demonstrates that the theoretical advantages
translate into concrete gains: $\protocol$ reduces the constraint count
by up to $7.5\times$ at $k = 1{,}024$ relative to a direct Freivalds SNARK
baseline, achieves up to $23.9\times$ proving-time speedup for
$k \geq 256$, and continues to scale at $k = 4{,}096$ where the baseline
exhausts memory.
The off-circuit encoding cost is negligible (under $1\,\text{s}$ at
$k = 2^{20}$), confirming that the code-based approach does not introduce
a new bottleneck.

\paragraph{Future directions.}
Several avenues remain.
First, while our protocol targets a single matrix multiplication, deep
learning inference involves \emph{chains} of matrix multiplications
interleaved with non-linear activations.
Extending $\protocol$ to handle layered computations---by chaining the
commitment outputs of one layer as inputs to the next---is a natural step
toward fully verifiable neural network inference.
Second, the encoding constraints, which contribute $\mathcal{O}(k^2)$ under
general constant-rate codes, could potentially be deferred to a separate
proximity test (e.g., FRI~\cite{ben2018fast}), further reducing the circuit.
Third, replacing the Merkle-based commitment with an algebraic vector
commitment (e.g., KZG) would reduce proof sizes from megabytes to kilobytes
at the cost of a structured setup; evaluating this trade-off in practice is
an important engineering task.
Finally, exploring \emph{batched verification}---where multiple matrix
multiplications share a single encoding and commitment phase---could yield
significant amortized cost savings in practical deployments.
